\documentclass[11pt,a4paper]{article}

\input{../../sablony/pripravy/priprava.sty}

\nastavPripravu{4. ročník}{Kombinatorika}{37}{Kombinace II -- podmínky a smíšené úlohy}

\begin{document}

\blok{Motivace a návaznost}

Při výběru komise často požadujeme určitý počet dívek a chlapců, přítomnost člena či zákaz společné účasti. Propojíme kombinace s permutacemi a variacemi.

\poznamka{Předpoklady: kombinace, variace, permutace a pravidlo součtu. Cíl: rozdělit výběr na části a samostatně rozhodnout o použité metodě.}

\blok{Výklad a odvození}

Při podmínce \uv{právě dva z jedné skupiny} použijeme součin kombinací za jednotlivé skupiny. Podmínku \uv{alespoň} často rozdělíme na disjunktní případy nebo spočteme doplněk. U smíšených úloh se před výpočtem vždy ptej, zda zaměnění pořadí osob vytváří nový výsledek.

\blok{Společně řešené příklady}

\begin{prikladReseniBox}[title={Příklad 1}]
Ze šesti dívek a pěti chlapců vybereme čtyřčlennou komisi s přesně dvěma dívkami. Kolik komisí vznikne?
\tcblower
$\binom62\binom52=15\cdot10=150$.
\end{prikladReseniBox}



\begin{prikladReseniBox}[title={Příklad 2}]
Z deseti lidí včetně Anny a Borise vybíráme čtyřčlenný tým. Kolik týmů obsahuje Annu, ale nikoliv Borise?
\tcblower
Anna je pevně vybrána, Boris zakázán, ostatní tři volíme z osmi: $\binom83=56$.
\end{prikladReseniBox}



\begin{prikladReseniBox}[title={Příklad 3}]
Ze sedmi lidí obsazujeme tři různé funkce. Kolik možností? A kolik je tříčlenných komisí bez funkcí?
\tcblower
Funkce $V(7,3)=210$, komise $\binom73=35$.
\end{prikladReseniBox}



\blok{Samostatná práce}

\begin{zadaniReseniBox}[title={Samostatná práce 1}]
Z pěti dívek a čtyř chlapců vybereme tři osoby, přesně dvě dívky.
\tcblower
$\binom52\binom41=40$.
\end{zadaniReseniBox}



\begin{zadaniReseniBox}[title={Samostatná práce 2}]
Ze šesti žáků vybereme dvojici, v níž nesmějí být společně Adam a Eva.
\tcblower
Všechny dvojice $\binom62=15$ minus zakázaná $1$: $14$.
\end{zadaniReseniBox}



\begin{zadaniReseniBox}[title={Samostatná práce 3}]
Pět různých knih chceme jen seřadit na polici. Jde o variace, permutace nebo kombinace?
\tcblower
O permutace, protože se řadí všech pět různých knih: $5!=120$.
\end{zadaniReseniBox}



\blok{Rozšíření}

\begin{zadaniReseniBox}[title={Rozšíření 1}]
Ze šesti dívek a pěti chlapců vybereme čtyřčlennou komisi obsahující alespoň dvě dívky.
\tcblower
Případy 2, 3, 4 dívky: $\binom62\binom52+\binom63\binom51+\binom64\binom50=150+100+15=265$.
\end{zadaniReseniBox}



\blok{Domácí úkol}

\poznamka{Doporučená práce doma: přibližně 15--30 minut. Úlohy jsou seřazeny od základních k náročnějším; učitel může poslední podúlohu označit jako volitelnou.}

\begin{zadaniReseniBox}[title={Domácí úkol}]

\begin{enumerate}[label=\alph*),leftmargin=*,itemsep=2pt]

\item Ze sedmi žen a pěti mužů vyber čtyřčlennou komisi s přesně dvěma ženami.
\item Z osmi osob (včetně A a B) vyber tříčlennou skupinu tak, aby A byl členem a B nebyl.
\item Kolik tříčlenných skupin lze vytvořit ze sedmi žáků, jestliže dva určení žáci nesmějí být společně?
\item Z osmi osob se zvolí předseda, zástupce a zapisovatel. Kolik různých obsazení?
\item Ze sedmi různých knih vybereme čtyři, potom je seřadíme. Kolik výsledků?

\end{enumerate}

\tcblower

\begin{enumerate}[label=\alph*),leftmargin=*,itemsep=2pt]

\item $\binom72\binom52=21\cdot10=210$.
\item Ze šesti ostatních vybereme dvě: $\binom62=15$.
\item $\binom73-\binom51=35-5=30$.
\item $V(8,3)=8\cdot7\cdot6=336$.
\item $\binom74\cdot4!=V(7,4)=840$.

\end{enumerate}

\end{zadaniReseniBox}

\poznamka{Orientační práce v hodině: 0--10 min srovnání; 10--25 min podmíněné kombinace; 25--40 min smíšené úlohy; 40--45 min vyhodnocení.}

\poznamka{Domácí úkol je mix tří typů bez opakování; zkontroluj hlavně zdůvodnění metody, ne pouze konečná čísla.}

\end{document}
